\chapter{Estat de l'art}
\label{cap:arte}

El fil d'aquest capítol és el mateix que el del problema (capítol~\ref{cap:intro}): un estudiant amb PDF propis no necessita un model de llenguatge aïllat, ni una app de flashcards aïllada, ni un cercador aïllat. Necessita un \textbf{cicle} amb quatre baules encadenades (figura~\ref{fig:cicle}):

\begin{enumerate}
  \item \textbf{Indexar} documents privats de manera recuperable (trossejar i representar).
  \item \textbf{Recuperar} el fragment adequat (i poder-lo citar).
  \item \textbf{Generar} una resposta o un material de pràctica \emph{ancorat} a aquest fragment.
  \item \textbf{Programar el repàs} perquè el que s'ha practicat no es perdi.
\end{enumerate}

\begin{figure}[htbp]
\centering
\begin{tikzpicture}[
  font=\small,
  box/.style={draw, rounded corners=2pt, align=center, inner sep=6pt,
    minimum width=2.55cm, minimum height=1.35cm},
  arr/.style={-{Latex[length=2.2mm]}, thick}
]
\node[box] (a) {1. Indexar\\PDF privat};
\node[box, right=0.45cm of a] (b) {2. Recuperar\\amb evidència};
\node[box, right=0.45cm of b] (c) {3. Generar\\ancorat};
\node[box, right=0.45cm of c] (d) {4. Repassar\\SM-2};
\draw[arr] (a) -- (b);
\draw[arr] (b) -- (c);
\draw[arr] (c) -- (d);
\end{tikzpicture}

\caption{Les quatre baules del cicle d'estudi que articulen aquest capítol.}
\label{fig:cicle}
\end{figure}

La literatura i el mercat cobreixen cada baula per separat, sovint amb molta profunditat. El que gairebé no existeix ---i el que aquest TFM reivindica--- és tancar el cicle de forma \emph{oberta, limitada i avaluable}, com un sistema d'enginyeria i no com un notebook. El detall de FastAPI, Next.js o Supabase no pertany aquí: és disseny (capítols~\ref{cap:arquitectura} i~\ref{cap:desarrollo}).

\section{De l'LLM nu al RAG}

Un LLM sense evidència externa respon amb el que ha memoritzat a l'entrenament. Això falla en apunts d'una assignatura concreta: el model no ha vist \emph{aquest} PDF, i si improvisa no hi ha manera d'auditar-lo. Lewis et al. \parencite{lewis2020rag} van formular la generació augmentada per recuperació (RAG) com la combinació d'un retriever i un generador condicionat als documents recuperats. La promesa és reduir al·lucinacions: la resposta ha de dependre d'un context extern, no només dels paràmetres del model.

A la pràctica, la qualitat depèn més de \emph{què} es recupera que de la mida de l'LLM \parencite{gao2024retrieval}. Per a aquest treball, això talla una via: no cal un model fronterer; cal un retriever honest, cites localitzables i un prompt que prohibeixi inventar fora del context. RAG cobreix, doncs, les baules~2 i~3 \emph{quan la tasca és preguntar}. Encara no diu com s'indexa un PDF de classe, ni com es practica després.

\section{Indexar: trossejar abans de recuperar}

La baula~1 no és ``pujar un fitxer a un xat''. És decidir \emph{quin} fragment entrarà a l'índex. Un PDF de 60 pàgines no cap sencer al retriever; es parteix en \emph{chunks}. Si el tall cau a mitja definició, ni BM25 ni l'embedding recuperen una unitat útil. Si els chunks són massa llargs, el model pot ``perdre'' el fet rellevant al mig del context \parencite{liu2024lost}. Barnett et al. \parencite{barnett2024seven} resumeixen set punts de fallada d'un RAG en producció; els primers són d'indexació (contingut absent o mal trossejat), no de l'LLM.

La taula~\ref{tab:chunking} contraposa estratègies habituals. Studyraft no inventa un algorisme nou: aplica una cascada (capçaleres $\rightarrow$ paràgrafs $\rightarrow$ oracions $\rightarrow$ tall dur) amb màxim 1200 caràcters i solapament 150 (capítol~\ref{cap:desarrollo}). El solapament redueix frases partides; duplica uns quants tokens d'embedding, un cost negligible (secció~\ref{sec:cost-api}).

\begin{table}[htbp]
\centering
\caption{Estratègies de trossejat (chunking) i el risc que arrosseguen.}
\label{tab:chunking}
\small
\begin{tabularx}{\textwidth}{lXX}
\toprule
Estratègia & Idea & Risc típic \\
\midrule
Finestra fixa & $N$ caràcters, independent de l'estructura. & Talla definicions i títols. \\
Finestra + solapament & El final d'un chunk es repeteix al següent. & Recupera ponts; duplica tokens. \\
Per capçaleres & Un chunk $\approx$ una secció del PDF. & Molts apunts no tenen títols detectables. \\
Studyraft (aquest TFM) & Cascada semàntica, $N{=}1200$, solapament 150, pàgines al metadada. & Sense OCR; un escaneig no genera chunks. \\
\bottomrule
\end{tabularx}
\end{table}

\section{Recuperar sobre apunts reals}

La baula~2 no és un únic algoritme. La cerca densa (embeddings + índex vectorial) captura paràfrasis: ``què és el potencial d'acció?'' \parencite{karpukhin2020dpr}. La cerca lèxica (BM25 / FTS) captura identificadors, fórmules i termes rars: ``què diu de $Na^+$ a la pàgina 12?'' \parencite{robertson2009probabilistic}. Un corpus d'apunts d'enginyeria barreja totes dues consultes; quedar-se només amb embeddings deixa escapar el que l'estudiant copia literalment del PDF.

Quan no es vol entrenar un model de fusió, Reciprocal Rank Fusion és un component estàndard per combinar llistes \parencite{cormack2009rrf}. Un rerank amb LLM millora sovint el top-$k$, a costa de latència i tokens: és un compromís de producte, no un requisit teòric (al capítol~\ref{cap:resultados} el p50 passa de 369\,ms a 1\,608\,ms). El fil és clar: RAG exigeix evidència; sobre apunts, l'evidència s'obté millor amb senyal densa \emph{i} lèxica, fusionades, i cites que mostrin fitxer i pàgina.

\begin{table}[htbp]
\centering
\caption{Senyals de recuperació sobre apunts (no és un ranking experimental).}
\label{tab:senyals}
\small
\begin{tabularx}{\textwidth}{lXX}
\toprule
Senyal & Què encerta & Què falla sol \\
\midrule
Dens (embedding) & Paràfrasi, sinònims, pregunta ``en les meves paraules''. & Sigles, fórmules, números de capítol. \\
Lèxic (BM25/FTS) & Termes rars i còpia literal del PDF. & Reformulacions sense paraules en comú. \\
Híbrid + RRF & Uneix les dues llistes sense un model de fusió entrenat. & Encara pot situar un chunk irrellevant al top-$k$. \\
+ rerank LLM & Reordena llegint pregunta i fragment. & Latència i cost; no substitueix un bon índex. \\
\bottomrule
\end{tabularx}
\end{table}

\section{Fidelitat: que la resposta estigui al chunk}

Recuperar un chunk no garanteix que la frase generada en depengui. La \emph{groundedness} (o fidelitat) és la propietat que cada afirmació de la resposta estigui suportada pel context recuperat. Protocols com RAGAS operacionalitzen això amb mètriques automàtiques (p.~ex. \emph{faithfulness}) \parencite{es2024ragas}. Gao et al. \parencite{gao2024retrieval} insisteixen que avaluar només la fluïdesa de l'LLM amaga al·lucinacions ben redactades.

Aquest TFM \textbf{no tanca} aquesta mètrica: exigiria anotació o un jutge LLM amb protocol (capítol~\ref{cap:resultados}). El que sí fa el sistema és \emph{reduir la superfície} d'invenció: prompt que prohibeix sortir del context, cites amb snippet, i (opcional) filtre lecture-focused. Això és enginyeria de producte, no una prova de fidelitat.

\section{Quan una sola passada no basta}

Fins aquí s'assumeix una pregunta puntual i uns quants chunks. Les preguntes d'estudi sovint són àmplies (``de què va aquest tema?'') o de capítol. Un índex jeràrquic (sinopsi de document $\rightarrow$ fragments) ajuda a no quedar-se només amb el fragment la frase del qual s'assembla a la pregunta. Els enfocaments agentic ---reescriure la consulta i tornar a recuperar, o crítiques com Self-RAG--- milloren la cobertura \parencite{asai2024selfrag}, però converteixen un \texttt{retrieve} en un bucle: cost, no-determinisme i més superfície d'error.

Per a un TFM, aquestes extensions existeixen, tenen sentit en un subconjunt de preguntes i s'han d'acotar (flags, una sola reescriptura). El cicle d'estudi no s'atura en recuperar millor: encara falta convertir l'evidència en pràctica.

\section{Més enllà del xat: pràctica i calendari}

Un xat citat tanca una sessió de dubtes; no consolida memòria. La baula~3, en un assistent d'estudi, també és \emph{generar} flashcards i qüestionaris a partir del mateix índex, i conservar-los com a conjunts editables. La baula~4 és programar el repàs.

SM-2 \parencite{wozniak1990supermemo} continua sent l'algoritme de referència pedagògic (Anki el va popularitzar). No és el model de memòria més recent: Half-Life Regression estima la vida mitjana d'un ítem a partir de l'historial de pràctica \parencite{settles2016hlr}; FSRS ajusta intervals amb més paràmetres a partir dels logs de ressenya \parencite{ye2024fsrs}. En un TFM d'enginyeria, la traçabilitat ---poder explicar per què una carta surt avui amb una fórmula tancada al domini--- pesa més que un petit guany de calendari. RAG dóna el context; SM-2 evita que el context sigui només una conversa efímera.

\section{Sistemes acadèmics (no només productes)}

La taula~\ref{tab:relacionados} compara apps de mercat. El tribunal també pregunta què hi ha a la \emph{literatura} de sistemes, no al catàleg comercial. Dues famílies són útils com a contrast, no com a rivals a ``guanyar'':

\begin{itemize}
  \item \textbf{RAG sobre documents amb cites.} PaperQA recupera fragments de papers científics i obliga a citar-los en la resposta \parencite{lala2023paperqa}. Comparteix amb Studyraft el problema de la baula~2+3 (evidència localitzable). No tanca flashcards ni SRS sobre els apunts de l'estudiant.
  \item \textbf{Tutoria intel·ligent (ITS).} Sistemes com AutoTutor i la revisió de VanLehn mostren que un tutor amb \emph{model de l'alumne} pot acostar-se, en alguns dominis, a un tutor humà \parencite{graesser2005autotutor,vanlehn2011tutoring}. Studyraft no modela l'estudiant més enllà de SM-2: no diagnostica concepcions errònies ni adapta un currículum.
\end{itemize}

\begin{table}[htbp]
\centering
\caption{Contrast acadèmic: cada sistema tanca un tros del cicle.}
\label{tab:academics}
\small
\begin{tabularx}{\textwidth}{lXX}
\toprule
Sistema & Què cobreix & Què no és aquest TFM \\
\midrule
PaperQA & RAG + cites sobre papers. & No és un cicle d'estudi (cards/SRS) sobre PDF de classe. \\
ITS (AutoTutor / VanLehn) & Tutoria amb model de l'alumne. & No indexa els apunts privats de l'estudiant. \\
HLR / Duolingo & SRS a escala amb model estadístic. & No recupera fragments d'un PDF propi. \\
Studyraft & Quatre baules, MVP avaluable. & Sense model cognitiu ric ni mètrica de fidelitat. \\
\bottomrule
\end{tabularx}
\end{table}

\section{Per què això ha de ser un sistema, no un notebook}

Totes les baules anteriors es poden demostrar en un colab: un PDF, uns embeddings, un prompt. Un TFM d'enginyeria informàtica ha d'argumentar \emph{modularitat i aïllament}: usuaris que no es veuen les dades, cues d'ingestió, substitució de proveïdor LLM, tests amb dobles. L'arquitectura hexagonal (ports i adaptadors) \parencite{cockburn2005hexagonal} és el marc clàssic per a això: el cas d'ús no coneix ni HTTP ni el client de base de dades.

Aquesta secció no compara frameworks de moda. Justifica una exigència del tribunal: el cicle d'estudi ha de ser llegible al codi i limitable (sense OCR, sense agent general, amb RLS). Les tecnologies concretes es trien al capítol~\ref{cap:arquitectura}.

\section{El mercat: cada producte tanca un tros del cicle}

Els productes comercials confirmen el diagnòstic. Anki i Quizlet excel·leixen a la pràctica i, en graus diversos, al calendari, però no parteixen dels PDF de l'estudiant amb cites. NotebookLM i el xat amb fitxers indexen documents i citen, però no tanquen flashcards, quiz persistents i SRS com a cicle d'estudi. La taula~\ref{tab:relacionados} és una \textbf{cobertura declarada} (agost de 2026, ús de les interfícies públiques conegudes), \emph{no} un ranking ni una avaluació empírica: Studyraft no «guanya»; ocupa la casella d'integració que aquest TFM construeix.

\begin{table}[htbp]
\centering
\caption{Cobertura declarada del cicle d'estudi (agost 2026). Pràctica = flashcards i qüestionaris. No és un ranking de qualitat.}
\label{tab:relacionados}
\small
\begin{tabular}{lcccc}
\toprule
Sistema & PDF propis & Cites & Pràctica & SRS \\
\midrule
Anki & no natiu & no & sí & sí \\
Quizlet & limitat & no & sí & limitat \\
NotebookLM & sí & sí & limitat & no \\
ChatGPT + fitxers & sí & variable & no & no \\
Studyraft & sí & sí & sí & sí (SM-2) \\
\bottomrule
\end{tabular}
\end{table}

\section{Posicionament}

El forat que ocupa aquest treball no és un nou algoritme de retrieval de primer nivell, ni un nou model de memòria. És la \textbf{integració avaluable} de les quatre baules ---indexar, recuperar amb cites, practicar, repassar--- sobre documents privats, amb una arquitectura que es pugui llegir, provar i limitar. Les seccions anteriors expliquen \emph{per què} cada baula existeix a la literatura; els capítols de disseny i desenvolupament expliquen \emph{com} Studyraft les encadena.
